\slajdid{10-s022}
\begin{frame}[label=10-s022]
Da bi kolo što manje opterećivalo neko prethodno kola ova struja bi trebalo da bude što je manje moguća. U tom smislu treba otpornik \underline{$R_B$ da bude što je moguće veći}. Međutim, prilikom analize karakteristike prenosa smo videli da treba otpornik \underline{$R_B$ da bude što je moguće manji} da bi imali što je veće pojačanje i užu prelaznu zonu. I ovakve situacije će nas često pratiti u elektronici. Da imamo suprotne zahteve prilikom izbora komponenti. Rešenje je u kompromisu. Nećemo se baviti u ovom trenutku kako napraviti kompromis, pošto je povezan i sa tačnošću, preciznošću, komponenti i sa disipacijom i sa … U analizi ćemo smatrati da je ta vrednost određena.
\end{frame}
